% ============================================================================
% METHODOLOGY
% ============================================================================

\section{Methodology}

\subsection{Problem Definition and Notation}

A Constraint Satisfaction Problem is formally defined as $\mathcal{P} = (\mathbf{X}, \mathbf{D}, \mathbf{C})$ where:
\begin{itemize}
    \item $\mathbf{X} = \{x_1, x_2, \ldots, x_n\}$ is a set of variables
    \item $\mathbf{D} = \{D_1, D_2, \ldots, D_n\}$ specifies finite domains ($x_i \in D_i$)
    \item $\mathcal{C} = \{c_1, c_2, \ldots, c_m\}$ is a set of constraints
\end{itemize}

A solution is a complete assignment $\theta: \mathbf{X} \to \bigcup_i D_i$ such that $\theta(x_i) \in D_i$ and $\forall c_j \in \mathcal{C}: c_j(\theta)$ holds.

\textbf{Solving Trajectory:} For puzzle instance $\mathcal{P}$, a solving trajectory is a temporal sequence of reasoning states:
$$\mathcal{T} = \langle (S_0, a_1, S_1), (S_1, a_2, S_2), \ldots, (S_{T-1}, a_T, S_T) \rangle$$

where $S_t = (C_t, \delta_t)$ is the solver state at step $t$, with $C_t$ being formalized constraints and $\delta_t$ being domain assignments. Action $a_t$ is LLM-generated inference (natural language statement + Z3 formalization). Trajectory $\mathcal{T}$ is \textit{successful} if $S_T$ satisfies all constraints with Z3 returning SAT and assignment matching ground truth.

\textbf{Solving Paradigm:} We define a solving paradigm as a five-tuple:
$$P = (\mathcal{Q}, \mathcal{O}, \phi_{\text{pre}}, \phi_{\text{post}}, \mathcal{S})$$

where:
\begin{itemize}
    \item $\mathcal{Q}$ specifies trigger conditions (constraint types and solver state features)
    \item $\mathcal{O}$ is the inference operation (parameterized template, e.g., $\texttt{pos}(X) = k + d$)
    \item $\phi_{\text{pre}}$ is the Z3-formalized precondition assertion
    \item $\phi_{\text{post}}$ is the Z3-formalized postcondition assertion
    \item $\mathcal{S}$ documents applicable scope (puzzle types and scales)
\end{itemize}

$\phi_{\text{pre}}$ and $\phi_{\text{post}}$ are written in SMT-LIB form so they can be mechanically checked by Z3 at each application—giving the paradigm a structured, executable interface that prior natural-language memory formats lack. We do not claim absolute formal correctness: the screening protocol described in Section~3.3.4 is sample-based and provides probabilistic acceptance guarantees rather than first-principles proofs.

\subsubsection{Paradigm: Position in the Inference-Rule Spectrum}
\label{sec:paradigm-spectrum}

To prevent confusion with adjacent notions, we briefly locate PRISM paradigms in the spectrum of learned inference rules:

\begin{itemize}
    \item \textbf{Lemma / learned clause (SAT/SMT):} Ground-truth, proof-derived from a single solver invocation, scoped to that invocation; correctness is automatic.
    \item \textbf{Macro-operator (EBL):} Generalised from one or more successful traces under a hand-crafted domain theory; provably sound by construction; scope is the domain theory's expressive range.
    \item \textbf{Unverified prior LLM-memory unit:} Natural-language heuristic mined from agent trajectories; cross-task reusable; quality uncertified — relies on subjective LLM self-evaluation.
    \item \textbf{PRISM paradigm (this work):} LLM-induced from clustered successful trajectories, then \emph{empirically} screened via Z3 sampling along three axes. It is cross-instance reusable (unlike a lemma), has no hand-built domain theory (unlike EBL), and carries an SMT-executable interface and a non-trivial empirical quality certificate (unlike prior LLM-memory units).
\end{itemize}

PRISM paradigms thus sit between provably-sound EBL macros (limited by the cost of a formal domain theory) and unverified LLM heuristics (limited by lack of mechanical quality control), trading a probabilistic acceptance guarantee for substantially wider applicability.

\subsection{PRISM Framework Overview}

PRISM operates via two decoupled stages: offline distillation and online inference.

\textbf{Offline Phase:} Given training puzzles $\{\mathcal{P}_1, \ldots, \mathcal{P}_N\}$:
\begin{enumerate}
    \item Collect successful solving trajectories via LLM+Z3
    \item Identify Key Decision Points (domain reductions $\geq$2 or non-trivial step types)
    \item Cluster KDPs via complete-linkage agglomerative clustering on constraint-type features
    \item Abstract paradigms via LLM synthesis from cluster representatives
    \item Filter candidates through triple Z3 verification (soundness $\geq$0.90, effect $\geq$0.80, precision $\geq$0.80)
    \item Persist verified paradigms in Paradigm Library $\mathcal{L}$
\end{enumerate}

\textbf{Online Phase:} For each puzzle $\mathcal{P}^*$:
\begin{enumerate}
    \item Extract constraint-type signature $\sigma(S_t) = \text{sorted}(\{type(c) \mid c \in C_t\})$
    \item Retrieve paradigms via Layer-1 (type matching) + Layer-2 (LLM semantic judgment)
    \item Execute Z3 consistency pre-check; inject passing paradigms as hints
    \item Generate and verify LLM inference steps via Z3
    \item Upon UNSAT, extract UNSAT core and perform causal attribution
    \item Record typed repairs in Repair Trajectory Memory $\mathcal{M}$
    \item Detect stagnation (Jaccard similarity) and loops (cosine similarity); trigger strategy escalation
    \item Write successful repairs back to augment $\mathcal{L}$
\end{enumerate}

\subsection{Offline Stage: Paradigm Distillation}

\subsubsection{Trajectory Collection}

Given $N = 600$ training puzzles across difficulty levels, we execute LLM+Z3 solving $r = 3$ times per puzzle with $\text{temperature} = 0.7$ for inference diversity. We retain only successful trajectories where Z3 returns SAT with assignments matching ground truth. This yields approximately 1500 successful trajectories as the distillation source.

Each trajectory step records: Z3 verification labels (SAT/UNSAT/ERROR), domain change records $\Delta_t$ (variables with reduced domains), constraint type signatures, and preliminary step classification.

\subsubsection{Key Decision Point Identification}

A step $(S_{t-1}, a_t, S_t)$ is a Key Decision Point if any of the following holds:
\begin{itemize}
    \item \textbf{Condition A (Significant Domain Reduction):} $\exists x: |\delta_{t-1}(x)| - |\delta_t(x)| \geq 2$ — at least one variable's domain shrinks decisively in a single step.
    \item \textbf{Condition B (Non-Trivial Step Type):} step classification is \textsc{Chain\_Propagation} or \textsc{Contradiction\_Elimination}.
    \item \textbf{Condition C (Aggregate Information Gain, soft):} the total log-ratio of domain reductions across all variables exceeds a bit-threshold $\tau_{\text{ig}}$:
    \[
    G(t) \;=\; \sum_{x: |\delta_t(x)| < |\delta_{t-1}(x)|} \log_2 \! \frac{|\delta_{t-1}(x)|}{|\delta_t(x)|} \;\geq\; \tau_{\text{ig}}\;\; (\text{default } 1.0 \text{ bit}).
    \]
\end{itemize}

Condition C complements the hard per-variable cutoff of A: it captures \emph{gradual tightening} steps that move several domains from $5\!\to\!3$ or $4\!\to\!2$ — individually below the threshold of A but collectively constituting a meaningful inference. Empirically the dominant useful pattern in harder puzzles (e.g.\ 6$\times$6 ZebraLogic) is precisely such gradual tightening, so condition C broadens paradigm coverage where the library most needs it. Step types are assigned via lightweight rule-based classification on natural-language action text.

\subsubsection{Cross-Trajectory Clustering and Paradigm Abstraction}

\textbf{Feature Representation:} Each KDP is represented as:
$$\mathbf{f}(\text{kdp}) = [\mathbf{h}_{\text{ct}}, \mathbf{b}_{\text{dom}}, \mathbf{e}_\tau, n_{\text{var}}, n_{\text{con}}]$$

where $\mathbf{h}_{\text{ct}}$ is a constraint-type histogram, $\mathbf{b}_{\text{dom}}$ is domain-size distribution, $\mathbf{e}_\tau$ is step-type one-hot encoding, and normalized problem scale metrics.

\textbf{Clustering:} We apply complete-linkage agglomerative clustering with cosine distance and threshold $\theta = 0.25$ (determined via grid search on development set). Clusters with fewer than 5 members are discarded.

\textbf{Paradigm Abstraction:} For each cluster, we uniformly sample up to 10 representative KDPs, format their solver states and reasoning text, and prompt an LLM to synthesize a reusable paradigm. LLM outputs are parsed into structured (trigger, operation, precondition, postcondition, scope) tuples. Up to 2 retries per cluster are permitted.

\subsubsection{Z3 Triple Verification}

For each candidate paradigm, we execute three independent verification checks:

\textbf{(1) Soundness:} Sample 50 random puzzle states satisfying trigger conditions. Test whether paradigm operations maintain Z3 satisfiability:
$$\text{soundness}(P) = \frac{1}{50} \sum_{i=1}^{50} \mathbf{1}[\text{Z3-SAT}(C_i \cup \{\text{op}(P)\})] \geq 0.90$$

\textbf{(2) Effect:} Verify that the paradigm operation is genuinely informative — not vacuous and not a tautology under the pre-condition.  Two sub-checks must both pass:
\begin{itemize}
    \item \textbf{Non-contradiction:} $\text{Z3-SAT}(\phi_{\text{pre}} \wedge \text{op}(P))$ — the operation is consistent with the pre-condition.
    \item \textbf{Non-vacuousness:} $\text{Z3-SAT}(\phi_{\text{pre}} \wedge \neg\,\text{op}(P))$ — there exists at least one state where the pre-condition holds but the operation does not. Equivalently, $\phi_{\text{pre}}$ does not entail $\text{op}(P)$. Were this check to fail (i.e.\ $\phi_{\text{pre}} \wedge \neg\,\text{op}(P)$ were UNSAT), the paradigm would be entailed by its own trigger and contribute no information.
\end{itemize}

\textbf{(3) Trigger Precision:} Estimate trigger selectivity against a representative constraint pool. Each pool entry is annotated with one or more canonical constraint kinds drawn from a fixed taxonomy $\mathcal{K} = \{$\texttt{position\_arithmetic}, \texttt{position\_equality}, \texttt{position\_inequality}, \texttt{attribute\_binding}, \texttt{distinct}, \texttt{quantified\_exclusion}, \texttt{integer\_bound}, \texttt{integer\_equality}, \texttt{integer\_inequality}, \texttt{integer\_arithmetic}$\}$. The trigger of paradigm $P$ fires on sample $s$ iff its declared constraint kinds $\mathcal{K}_P$ have non-empty intersection with the kinds present in $s$ —
a structural set-intersection match rather than substring containment, which prevents spurious lexical co-occurrence (e.g., a \texttt{position\_arithmetic} paradigm does not falsely fire on a generic attribute binding that mentions ``position'').
\[
\text{precision}(P) =
\frac{1}{N_p} \sum_{i=1}^{N_p} \mathbf{1}[\,\mathcal{K}_P \cap \mathcal{K}_{s_i} = \emptyset\,]
\geq \tau_p
\]

Here, $N_p = 50$ random subsets of size 2--5 are drawn from a pool of $|\text{pool}| = 18$ annotated constraints spanning ZebraLogic-domain and generic arithmetic patterns; $\tau_p$ is the non-firing rate floor.

Paradigms passing all three checks are persisted in Paradigm Library $\mathcal{L}$. Across 600 training puzzles, offline distillation yields approximately 35 verified paradigms (40:1 compression ratio).

\subsection{Online Stage: Paradigm-Guided Inference}

\subsubsection{Two-Layer Paradigm Retrieval}

For puzzle $\mathcal{P}^*$, at each step $t$, we extract constraint-type signature $\sigma(S_t)$ and apply:

\textbf{Layer-1 (Structured Type Matching):} Return paradigms $P \in \mathcal{L}$ ranked by a weighted combination of scope-Jaccard similarity with $\sigma(S_t)$ and confidence. Crucially, type-set inclusion alone is too weak: two puzzles with the same set of constraint types may differ substantially in multiplicity (e.g.\ one ordering constraint vs.\ four). To address this, each paradigm's trigger may carry optional \emph{relational predicates} from a small vocabulary --- \texttt{count\_atleast}($t$,$n$): at least $n$ constraints of type $t$ must be present; \texttt{count\_atmost}($t$,$n$); \texttt{cooccur}($t_1$,$t_2$): both types must be present. Predicates are evaluated against the actual type bag (a list with repetitions) and paradigms whose predicates fail are pruned \emph{before} ranking. This lifts Layer-1 from pure set-membership to a set-with-multiplicity match while remaining LLM-free and $O(|\mathcal{L}|)$. The combined filter typically reduces candidates from 35 to $< 5$.

\textbf{Layer-2 (Semantic Filtering, cost-aware):} Layer-2 invokes an additional LLM call to resolve ambiguity among Layer-1 candidates. Because this is the most expensive per-step component, we gate it by an activation policy. Under our default \texttt{complexity\_gated} policy, Layer-2 runs only when (i) Layer-1 returns $\geq 2$ candidates \emph{and} (ii) either the puzzle size meets a complexity floor (default 25 constraints, $\approx 5\!\times\!5$ ZebraLogic) or at least one UNSAT has been observed in the current puzzle. Alternative policies \texttt{always} (legacy behaviour) and \texttt{stagnation\_only} (Layer-2 deferred until the repair memory flags stagnation) are also supported. When Layer-2 is gated off, the Layer-1 top-1 candidate is retained — guidance is never fully suppressed, and the downstream Z3 consistency pre-check still filters obviously incompatible hints.

\subsubsection{Consistency Pre-Check and Hint Injection}

Before paradigm injection, execute Z3 consistency verification:
$$\text{consistent}(P, S_t) = \text{Z3-SAT}(C_t \cup \{\text{op}(P)\})$$

Only paradigms returning SAT are injected as "soft" (advisory) hints, not hard constraints. This prevents LLM misapplication when paradigms don't fully align with current state.

\subsubsection{Real-Time Z3 Verification and UNSAT Attribution}

After LLM generates inference action $a_t$, add to solver via constraint registry (enabling tracking) and call Z3:

\textbf{SAT Case:} Update domain assignments, save as potential revert checkpoint. Continue solving.

\textbf{UNSAT Case:} Extract UNSAT core $U_t$ and perform two-level causal attribution.

\textbf{Primary attribution} (coarse):
$$\text{attr}(a_t, U_t) = \begin{cases} \text{NEW\_ASSERTION} & \text{if } a_t \in U_t \\ \text{LEGACY\_ERROR} & \text{if } a_t \notin U_t \end{cases}$$

If \textsc{Legacy\_Error}, the contradiction predates $a_t$; revert $a_t$ and diagnose the pre-existing set.  If \textsc{New\_Assertion}, $a_t$ itself is responsible.

\textbf{Secondary classification} (fine-grained) refines \textsc{New\_Assertion} into a structured \textit{ErrorType}. Classification is performed by an AST pass over the parsed Z3 expression rather than by substring matching on the raw constraint string — this makes the classifier robust to Z3's many equivalent surface forms (e.g.\ \texttt{Not(x > 0)} is recognised as a flipped form of \texttt{x <= 0}).
\begin{itemize}
    \item \textbf{SyntaxError} — the assertion fails Z3 parsing.
    \item \textbf{ScopeError} — the parsed expression references an identifier outside the currently-declared variable set.
    \item \textbf{SemanticFlip} — the AST contains contradictory comparison operators on identical operand pairs (e.g.\ $x \geq k \wedge x \leq k\!-\!1$, or $x = k \wedge \neg(x = k)$). Detection traverses the expression with a polarity-tracking walk so that negations under \texttt{Not} are normalised before atom comparison.
    \item \textbf{OverConstraint} — default: the new assertion legally tightens the formula beyond satisfiability.
\end{itemize}

The \textit{ErrorType} label is stored in $R_t$ and used by the strategy escalation to select the most appropriate repair class (e.g., SemanticFlip triggers constraint-direction repair rather than bound relaxation).

\subsection{Repair Trajectory Memory}

\subsubsection{Memory Data Structure}

Repair Trajectory Memory $\mathcal{M}$ maintains ordered records $\{R_1, \ldots, R_L\}$, each a six-tuple:
$$R_t = (e_t, U_t, \hat{U}_t, \alpha_t, o_t, U'_t)$$

where:
\begin{itemize}
    \item $e_t$ is error type classification
    \item $U_t$ is UNSAT core (constraint IDs)
    \item $\hat{U}_t = \text{SHA256}(\text{sorted}(U_t))$ is core fingerprint for rapid comparison
    \item $\alpha_t$ is repair action (type, target constraint, summary text, embedding)
    \item $o_t$ is repair outcome (SAT or UNSAT)
    \item $U'_t$ is new UNSAT core post-repair (if UNSAT)
\end{itemize}

Repair action embeddings are computed via all-MiniLM-L6-v2 sentence transformer (384-dim), enabling semantic similarity queries.

\subsubsection{Stagnation Detection}

Stagnation occurs when recent repair records show high UNSAT core overlap, indicating that repairs are not addressing the root contradiction:
$$\text{stagnate}(\mathcal{M}, k=3) = \mathbf{1}\!\left[\max_{i \neq j \in \text{last-}k} J(\tilde{U}_i, \tilde{U}_j) \geq \tau_s\right]$$

where Jaccard similarity $J(A,B) = |A \cap B| / |A \cup B|$ is computed over \emph{canonicalised} cores $\tilde{U}_t$ (see below) drawn from the last $k = 3$ records.  We use the \emph{maximum} (any-pair) rather than the minimum (all-pairs) as the trigger criterion: early-warning detection is preferable to late-stage detection because the former preserves more of the repair budget for recovery actions.

\paragraph{Z3 core non-uniqueness and canonicalisation.} Z3's \texttt{unsat\_core()} returns \emph{one} minimal unsatisfiable subset, not the unique one — different solver invocations on the same logically unsatisfiable constraint set may return different cores depending on assumption ordering and internal heuristics. Naively comparing raw cores can therefore both \emph{under-detect} stagnation (the same underlying contradiction surfaces as different core sets across iterations) and \emph{over-detect} it (two unrelated contradictions share a common high-frequency constraint such as \texttt{Distinct(...)}). To mitigate this we canonicalise each core via (i) constraint-string normalisation (whitespace, operator aliasing such as \texttt{Not(x>0)} $\leftrightarrow$ \texttt{x<=0}), and (ii) variable renaming to a canonical schema (\texttt{v1, v2, ...}) by first-occurrence order. The fingerprint $\hat{U}_t$ stored in each repair record is the SHA-256 of the canonicalised core, making cross-iteration comparison robust to syntactically-equivalent but lexically-different core listings.

\subsubsection{Loop Detection}

Loop detection identifies when a proposed repair duplicates a historical attempt. We combine two signals so that LLM paraphrase noise does not defeat detection:

\textbf{(i) Structured signal (primary).} Each repair action carries a structured triple $(\tau, c, \pi)$ where $\tau$ is the repair type (drawn from a closed vocabulary of eight classes such as \texttt{relax\_bound}, \texttt{split\_constraint}), $c$ is the target constraint id, and $\pi$ is a parameter signature (e.g.\ \texttt{+1}, \texttt{*2}, \texttt{removed\_clause\_3}) summarising the atomic edit. A loop is flagged when this triple exactly matches that of any prior record, regardless of natural-language phrasing.

\textbf{(ii) Semantic signal (fallback).} When the structured triple is absent or does not match, we fall back to cosine similarity between sentence-transformer embeddings of the action summaries:
$$\text{loop}(\alpha_{\text{new}}) = \mathbf{1}\!\left[\,\big((\tau, c, \pi)_{\text{new}} = (\tau, c, \pi)_t \text{ for some } t\big)\,\vee\,\big(\max_t \cos(\mathbf{e}_{\text{new}}, \mathbf{e}_t) \geq \tau_\ell\big)\right]$$
with $\tau_\ell = 0.90$. The structured-signal-first design lets PRISM detect loops even when the LLM paraphrases (which routinely produces cosines in the $0.78\text{-}0.85$ range — below the safe semantic threshold).

\subsubsection{Four-Level Strategy Escalation}

Upon stagnation or loop detection:
\begin{enumerate}
    \item \textbf{L1 - Switch Target:} Preserve repair type but target different constraint from UNSAT core
    \item \textbf{L2 - Switch Type:} Change repair strategy class (e.g., relax-bound $\to$ split-constraint)
    \item \textbf{L3 - Revert Checkpoint (multi-level stack):} Pop the most recent SAT-adjacent state from a bounded checkpoint stack (depth $\leq 5$) and roll back the solver to it; resume from an alternative inference path. Successive L3 escalations consume successive entries from the stack, enabling the system to roll back further when a single-step revert does not break stagnation. When the stack is exhausted, the next escalation falls through to L4.
    \item \textbf{L4 - Full Retranslation:} Discard current formalization; re-translate all constraints from natural language with error history as guidance
\end{enumerate}

\paragraph{Termination property.} The online phase is guaranteed to terminate within a bounded number of LLM invocations per puzzle. Concretely: (i) the outer repair loop is capped by \texttt{max\_repair\_rounds} ($R{=}5$ in our default configuration); (ii) L4 (Full Retranslation) is triggered at most once per puzzle — a second attempt at L4 is suppressed and the puzzle is marked unsolved; (iii) L3 (Revert Checkpoint) pops at most $|$checkpoint-stack$|$ entries, and an exhausted stack falls through to L4. Thus the worst-case per-puzzle LLM-call budget is $1\,(\text{translate}) + R\,(\text{repair rounds}) + 1\,(\text{L4 retranslate}) = R+2$. A short proof sketch is given in Appendix B.

\subsection{Paradigm Write-Back and Library Maintenance}
\label{sec:writeback}

When a repair (or a fresh inference step) succeeds, PRISM does \emph{not} immediately admit the corresponding pattern into the live paradigm library $\mathcal{L}$. We adopt a \emph{staged write-back} policy:

\begin{enumerate}
    \item \textbf{Candidate staging (online, $O(1)$).} On a successful repair / SAT continuation, the responsible (trigger features, operation, inferred pre/post-condition) tuple is appended to a candidate pool $\mathcal{L}^{\dagger}$. No verification is performed online, so the inference loop remains low-latency.
    \item \textbf{Batched re-verification (offline / periodic).} Every $K$ solved puzzles (or end-of-evaluation), candidates in $\mathcal{L}^{\dagger}$ are de-duplicated against $\mathcal{L}$ (by trigger-kind signature and operation-template hash) and re-run through the same triple verification protocol (Section 3.3.4) used in offline distillation.
    \item \textbf{Promotion criterion.} Candidates that pass all three checks (soundness $\geq \tau_s$, effect, precision $\geq \tau_p$) are admitted to $\mathcal{L}$ with their support count and confidence initialised from the screening trials; failed candidates are discarded with provenance retained for later analysis.
    \item \textbf{Existing-paradigm updates.} When an online step matches an existing $P \in \mathcal{L}$ and succeeds, only its $\langle\text{support\_count}, \text{confidence}\rangle$ counters are updated atomically — no schema change, no re-verification.
\end{enumerate}

This staged policy keeps the online critical path free of expensive verification while ensuring that $\mathcal{L}$ retains the same SMT-screened quality guarantee at all times. The candidate pool $\mathcal{L}^{\dagger}$ is not used for retrieval during the online phase.

\subsection{Computational Cost Analysis}
\label{sec:cost}

PRISM's cost has two components that should be reported separately because they amortise differently across deployments.

\paragraph{Offline (one-time, amortised over the test set).}
The offline phase issues LLM calls in three places: (i) trajectory collection, $N_{\text{train}} \cdot r \cdot \bar{T}$ where $r{=}3$ is sampling repetition and $\bar{T}$ is the mean trajectory length; (ii) paradigm abstraction, at most $2|\mathcal{C}|$ calls for $|\mathcal{C}|$ clusters with up to two retries each; (iii) batched re-verification of online write-back candidates, $\propto |\mathcal{L}^{\dagger}|$. Z3 work scales as $50 \cdot |\mathcal{P}_{\text{candidate}}|$ verification trials. The entire offline pipeline runs once and its cost is amortised over every subsequent test puzzle: $C_{\text{off,amortised}} = C_{\text{offline,total}} / N_{\text{test}}$.

\paragraph{Online (per puzzle, worst case).}
Per the termination property (Section 3.5), the online phase invokes the LLM at most $R + 2$ times per puzzle: one translation, at most $R$ repair-round inferences, and at most one L4 retranslation. The Layer-2 semantic match adds an extra LLM call per step \emph{only} when the cost-aware policy fires (Section 3.4.1); in the default \texttt{complexity\_gated} policy this contributes $\leq R$ additional calls and is suppressed entirely on easy instances. Z3 work per step is bounded by one consistency pre-check + one solver call + one UNSAT-core extraction.

\paragraph{Amortised total per test puzzle.}
$$
C_{\text{total}}(p) \;=\; \underbrace{C_{\text{off,amortised}}}_{\text{shared offline cost}} \;+\; \underbrace{C_{\text{online}}(p)}_{\le (R{+}2) + \mathbb{1}[\text{L2 fires}] \cdot R \text{ LLM calls}}.
$$
For deployments where $N_{\text{test}}$ is large (typical), $C_{\text{off,amortised}}$ becomes negligible and the dominant cost is the bounded online budget. We report both components separately in Section~4 so that the comparison with neural-symbolic baselines reflects the realised, not headline, cost.
